\section{含辅助“重夸克”场的有效 QCD}

具体地，考虑如下非局域算符
\begin{equation}\label{Eq:nonlocalop}
 O(x,y) = \overline{\psi}(x) \Gamma L(x,y)\psi(y),
\end{equation}
其中 $\psi, \bar\psi$ 表示裸夸克场，$\Gamma$ 为某个狄拉克结构。$L(x, y)$ 是从 $y$ 到 $x$ 的按路径排序的规范链，
\begin{equation}
             L(x, y) = {\cal P}\exp\left(-ig\int^1_0 d\lambda\,\frac{dz^\mu}{d\lambda}  A_\mu(z(\lambda))\right)
\end{equation}
携带色 SU(3) 群的 $(3,\bar 3)$ 表示，其中路径的参数化满足 $z(0)=y$ 与 $z(1)=x$。该规范链保证算符 $O(x,y)$ 的规范不变性。
对于本文主要关心的直线规范链，我们可以选取 $z^\mu(\lambda)=y^\mu+\eta\lambda\, n^\mu$，其中 $n^\mu$ 为指定方向的单位矢量，$\eta$ 为规范链的长度。

为研究上述算符的重正化性质，我们引入一个“重夸克”辅助场（没有自旋自由度的色源）$Q$，色表示为 3，及其共轭 ${\overline Q}$，色表示为 ${\bar 3}$。
我们把 QCD 扩展以包含这种“重夸克”与胶子场的相互作用，并引入如下拉氏量
\beq\label{Eq:lag}
{\cal L}={\cal L_{\text{QCD}}}+\overline Q(x)i n\cdot D Q(x),
\eeq
其中 $D_\mu=\partial_\mu+igt^a A^a_\mu$ 是基础表示下的协变导数。对于真实重夸克或类时威尔逊线，$n^\mu$ 取类时矢量 $n^\mu=(1,0,0,0)$；而对于类空威尔逊线，$n^\mu$ 可以取为 $n^\mu=(0,0,0,1)$。下面我们将聚焦于后一情形，尽管该讨论对任意 $n$ 或许都适用。
只要时间分量 $n^0$ 非零，“重夸克”就是动力学的，但沿一条指定的世界线运动。

在上述理论中，我们可以用一个新算符替换双局域算符 $O(x,y)$：
\begin{equation}\label{Eq:repl}
             O(x, y) = \overline{\psi}(x)\Gamma Q(x) \overline{Q}(y) \psi(y).
\end{equation}
这可以如下看出。在对“重夸克”场积分之后，我们得到
\beq
\int {\cal D}\overline Q{\cal D}Q\, Q(x)\overline Q(y) e^{i\int d^4x {\cal L}}=S_Q(x, y) e^{i\int d^4x {\cal L_{\text{QCD}}}}
\eeq
相差一个行列式 ${\rm det}(in\cdot D)$，可以证明它是常数并可吸收进生成泛函的归一化之中~\cite{Mannel:1991mc}，其中 $S_Q(x, y)$ 满足
\beq
n\cdot D\, S_Q(x, y)=\delta^{(4)}(x-y),
\eeq
在 $n^\mu=(0,0,0,1)$ 的情形下，取适当的边界条件，其解为
\begin{align}
&S_Q(x,y)=\theta(x^z-y^z)\delta(x^0-y^0)\delta^{(2)}(\vec x_\perp-\vec y_\perp)L(x,y)\non\\
&=\theta(x^z-y^z)\delta(x^0-y^0)\delta^{(2)}(\vec x_\perp-\vec y_\perp)L(x^z,y^z),
\end{align}
以下不失一般性，我们将限于 $x^z>y^z$ 的情形。$\delta$ 函数保证 $x$ 与 $y$ 的时间分量和横向分量相等，从而沿纵向方向生成一条类空威尔逊线。上述结果使我们能够把准 PDF 中出现的威尔逊线 $L(x^z,y^z)$ 替换为两个辅助重夸克场之积 $Q(x^z)\overline Q(y^z)$。于是，准 PDF 中的非局域算符便成为式~(\ref{Eq:repl}) 中两个复合算符的乘积。
